\documentclass[a4paper]{book}

%%%%%%%%%%%%%%%%%%%%%%%%%%%%%%%%%%%%%%%%%%%%%%%%%%
% localization	
\usepackage[latin2]{inputenc}
\usepackage[magyar]{babel}

%%%%%%%%%%%%%%%%%%%%%%%%%%%%%%%%%%%%%%%%%%%%%%%%%%
% magyarul �rjuk a dokumentumot, �gy 
% magyar nyelv� komenteket haszn�llunk
\usepackage{amsmath}
\usepackage{amssymb}

\newcommand{\akkor}{\Rightarrow}
%%%%%%%%%%%%%%%%%%%%%%%%%%%%%%%%%%%%%%%%%%%%%%%%%%
% amsthm nem kompatibilis \usepackage[magyar]{babel}
% ez�rt _ntheorem_ csomagot hasznallunk.
\usepackage[amsmath,thmmarks,thref]{ntheorem}
\numberwithin{equation}{chapter}

%%%%%%%%%%%%%%%%%%%%%%%%%%%%%%%%%%%%%%%%%%%%%%%%%%
% _T�telek_ 				- thm 	k�rnyezet
% _�ll�t�sok_ 			- all 	k�rnyezet
% _Lemm�k_					- lem		k�rnyezet
% _Defin�ci�k_ 			- defn 	k�rnyezet
% _K�vetkezm�nyek_ 	- kov 	k�rnyezet
\theoremstyle{change}
%\theoremheaderfont{\bfseries}
\theorembodyfont{\normalfont}

\newtheorem{thm}{T�tel}[section]				% thm k�rnyezet
\newtheorem{all}[thm]{�ll�t�s}			% all k�rnyezet
\newtheorem{lem}[thm]{Lemma}				% lem k�rnyezet
\newtheorem{defn}{Defin�ci�}		% defn k�rnyezet
\newtheorem{kov}[thm]{K�vetkezm�ny}	% kov k�rnyezet

%%%%%%%%%%%%%%%%%%%%%%%%%%%%%%%%%%%%%%%%%%%%%%%%%%
% _Megjegyz�sek_ 	- meg k�rnyezet
% _P�ld�k_ -			- pld k�rnyezet
\theoremstyle{nonumberplain}
%\theoremheaderfont{\bfseries}
\theorembodyfont{\normalfont}

\newtheorem{meg}{Megjegyz�sek:} % meg k�rnyezet

\theoremstyle{nonumberbreak}
\newtheorem{pld}{P�lda:}
%%%%%%%%%%%%%%%%%%%%%%%%%%%%%%%%%%%%%%%%%%%%%%%%%%
% _Bizony�t�sok_ - biz k�rnyezet
\theoremstyle{nonumberplain}
\theoremheaderfont{\normalfont\itshape}
\theorembodyfont{\normalfont}
\theoremsymbol{\ensuremath{_\Box}} 

\newtheorem{biz}{Bizony�t�s.} % biz k�rnyezet

%%%%%%%%%%%%%%%%%%%%%%%%%%%%%%%%%%%%%%%%%%%%%%%%%%
% parancsok
% matematikai parancsok

\newcommand{\vt}[1]{\mathbf{\underline{#1}}}

%%%%%%%%%%
% halmazok
\newcommand{\field}[1]{\mathbb{#1}}
\newcommand{\N}{\field{N}}
\newcommand{\Z}{\field{Z}}
\newcommand{\Q}{\field{Q}}
\newcommand{\R}{\field{R}}
\newcommand{\K}{\field{K}}

%%%%%%%%%%
% rel�ci�k
\renewcommand{\div}{|}
\newcommand{\divu}{\|}
\newcommand{\ndiv}{\nmid}

\newcommand{\eq}{\backsim}
\newcommand{\nequiv}{\not\equiv}

%%%%%%%%%%
% abszolut ertek, norma
\providecommand{\abs}[1]{\left\vert #1\right\vert}
\providecommand{\norm}[1]{\left\Vert #1\right\Vert}

\newcommand{\normp}[2][p]{\abs{#2}_#1}

\newcommand{\lnko}[1]{\left({#1}\right)}
\newcommand{\gen}[1]{\left\langle  #1\right\rangle}

%%%%%%%%%%
% fuggvenyek
\DeclareMathOperator{\ord}{ord}
\newcommand{\ordp}[2][p]{\ord_{#1}\left(#2\right)}
\DeclareMathOperator{\met}{\rho}

\begin{document}
\begin{titlepage}
\begin{flushright}
E�tv�s Lor�nd Tudom�nyegyetem\\
Informatikai Kar \\ 		
Programoz�si Nyelvek �s Ford�t�programok Tansz�k
\end{flushright}
\includegraphics[bb = 0cm 0cm 0cm -2cm]{elte3.png} 

%\begin{tabular}{p{3cm}r}
%&  \small E�tv�s Lor�nd Tudom�nyegyetem\\
%& \small Informatikai Kar \\ 
%& \small Programoz�si Nyelvek �s Ford�t�programok Tansz�k
%\end{tabular}
\vspace{3cm}
\begin{center}\LARGE
GTL \\ Graphical Template Library
\par\end{center}
\begin{center} \large \lineskip 0.5em \begin{tabular}[t]{c}
{\bfseries Vatai Emil}\\
{\normalsize V. �ves Programtervez� Matematikus hallgat�}\\
\vspace{1cm}
{\normalsize T�mavezet�:}\\
Dr. {\bf Porkol�b Zolt�n} egyetemi docens\\
{\normalsize E�tv�s Lor�nd Tudom�nyegyetem - Informatikai Kar}\\
{\normalsize Programoz�si Nyelvek �s Ford�t�programok Tansz�k}

\end{tabular}\par\end{center}


\vspace{5cm}
Budapest, \today

\end{titlepage}

\end{document}


